% =====================================================================
%  Curriculum Vitae - Giyoon Jang (장지윤)
%  Compile with:  xelatex Curriculum_Vitae_GiyoonJang.tex
%  (XeLaTeX is required because the document contains Korean text.)
% =====================================================================

\documentclass[10pt,a4paper]{article}

% --- Korean language support -----------------------------------------
% Uses the kotex package when it is available (e.g. Overleaf, TeX Live full).
% Otherwise it falls back to a CJK-capable system font via fontspec.
\usepackage{iftex}
\ifXeTeX
    \usepackage{fontspec}
    \IfFileExists{kotex.sty}
        {\usepackage{kotex}}
        {\setmainfont{Noto Serif CJK KR}}
\else
    \usepackage[T1]{fontenc}
    \usepackage[utf8]{inputenc}
    \usepackage{kotex}
\fi
% ---------------------------------------------------------------------

\usepackage{geometry}
\geometry{left=1in, right=1in, top=1in, bottom=1in}
\usepackage{enumitem}
\usepackage{hyperref}
\usepackage{titlesec}
\usepackage{xcolor}
\usepackage{graphicx}

% Title formatting based on the original CV.
\titleformat{\section}{\large\bfseries}{}{0em}{}[\titlerule]

% Reusable CV commands.
\newcommand{\seclabel}[1]{\section*{\uppercase{#1}}}
\newcommand{\cvitem}[2]{\textbf{#1} \hfill \small{#2}\\}
\newcommand{\placeholder}[1]{\textnormal{[#1]}}
\hypersetup{
    colorlinks=true,
    linkcolor=cyan,
    filecolor=cyan,
    urlcolor=cyan,
    citecolor=cyan
}

\setlist[itemize]{leftmargin=0pt}
\setlist[enumerate]{leftmargin=0pt}

\begin{document}

\noindent
\begin{minipage}[c]{\textwidth}
    \begin{center}
        \textbf{\Huge{Giyoon Jang}} \\ \vspace{1.0em}
        M.S. Student, School of Software \\
        Soongsil University \\
        \href{https://sites.google.com/view/ssu-nlp/home}{Natural Language Processing Lab} \\
        \textbf{Email: }\href{mailto:jyjang1130@soongsil.ac.kr}{jyjang1130@soongsil.ac.kr} \\
        \textbf{Home: }\href{https://jiyoon113.github.io}{https://jiyoon113.github.io} \\
        \textbf{GitHub: }\href{https://github.com/jiyoon113}{https://github.com/jiyoon113} \\
        \textbf{LinkedIn: }\href{https://www.linkedin.com/in/giyoon-jang-784876165/}{https://www.linkedin.com/in/giyoon-jang-784876165/} \\
        \textbf{Lab: }\href{https://sites.google.com/view/ssu-nlp/home}{https://sites.google.com/view/ssu-nlp/home}
    \end{center}
\end{minipage}

\vspace{1em}

\seclabel{Lifelong Research Objective}
\small{My long-term objective is to make generative AI systems that are not only capable but also
trustworthy and genuinely useful to people, by studying how large language models should be
evaluated, governed, and safely deployed in real-world settings.}

\seclabel{Research Interests}
\small{Natural Language Processing, Large Language Models, Generative AI, LLM Evaluation and
Benchmarking, Trustworthy and Safe AI, Human-Centered AI Applications}

\seclabel{Education}
\begin{itemize}
    \item \cvitem{Soongsil University}{Mar 2026 -- Present}
    \small{M.S. in School of Software (advisor: \href{https://parkchanjun.github.io/}{Chanjun Park}) \\
    \href{https://sites.google.com/view/ssu-nlp/home}{Natural Language Processing Lab}} \vspace{0.6em}

    \item \cvitem{Hanyang University (ERICA, Ansan)}{Mar 2019 -- Aug 2025}
    \small{B.S. in Computer Science \\
    GPA: 3.82 / 4.5} \vspace{0.6em}

    \item \cvitem{University of California, San Diego (UCSD)}{Sep 2023 -- Jun 2024}
    \small{Exchange Student Program \\
    Coursework: Discrete Mathematics, Computational Theory}
\end{itemize}

\seclabel{Work Experience}
\begin{itemize}
    \item \cvitem{Soongsil University}{Mar 2025 -- Present}
    \small{Graduate Researcher \\
    \href{https://sites.google.com/view/ssu-nlp/home}{Natural Language Processing Lab}} \vspace{0.6em}

    \item \cvitem{Lotte Innovate}{Dec 2024 -- Feb 2025}
    \small{Intern \\
    Conducted AI image generation research and designed the back-end database structure using Firebase.} \vspace{0.6em}

    \item \cvitem{Einsis INC}{Jun 2022 -- Dec 2022}
    \small{IT Service Manager \\
    Streamlined cross-team processes and implemented workflow solutions for service operations.}
\end{itemize}

% --- PUBLICATIONS -----------------------------------------------------
% No publications yet. To add them later, uncomment this block by removing
% the leading "% " from each line.
% \seclabel{Publications}
% \textit{Equal contribution}: *, \textit{Corresponding author}: $\dagger$
% \vspace{0.4em} \\
% \noindent\textit{Google Scholar}: \href{https://scholar.google.com}{\placeholder{YOUR GOOGLE SCHOLAR URL}}
%
% \noindent\textit{Semantic Scholar}: \href{https://www.semanticscholar.org}{\placeholder{YOUR SEMANTIC SCHOLAR URL}}
%
% \noindent\textit{arXiv}: \href{https://arxiv.org}{\placeholder{YOUR ARXIV AUTHOR URL}}
%
% \paragraph{Conferences}
% \begin{enumerate}
%     \item \small{\textbf{\placeholder{PAPER TITLE}} \\
%     \placeholder{AUTHOR 1}, \underline{Giyoon Jang}\textsuperscript{$\dagger$}, \placeholder{AUTHOR 3} \\
%     \textit{\placeholder{CONFERENCE NAME, YEAR}}}
% \end{enumerate}
%
% \paragraph{Workshops}
% \begin{enumerate}
%     \item \small{\textbf{\placeholder{WORKSHOP PAPER TITLE}} \\
%     \placeholder{AUTHOR LIST WITH YOUR NAME UNDERLINED} \\
%     \textit{\placeholder{WORKSHOP NAME, YEAR}}}
% \end{enumerate}
%
% \paragraph{Journals}
% \begin{enumerate}
%     \item \small{\textbf{\placeholder{JOURNAL ARTICLE TITLE}} \\
%     \placeholder{AUTHOR LIST WITH YOUR NAME UNDERLINED} \\
%     \textit{\placeholder{JOURNAL NAME, VOLUME, ISSUE, PAGES, YEAR}}}
% \end{enumerate}

\newpage

\seclabel{Research Projects}
\textbf{Soongsil University}
\begin{itemize}[left=0pt]
    \item \small 2026년 인공지능 말평 운영 및 연계 행사 개최 (2026, 참여연구자 (용역과제), National Institute of the Korean Language)
    \item \small 2026년 인공지능의 한국어 능력 표준 평가 과제 구축 및 활용 (2026, 참여연구자 (용역과제), National Institute of the Korean Language)
    \item \small 우수연구 신진연구 (유형B) (2026-2029, 참여연구자 (3책5공), NRF)
    \item \small Foundation Model 운용과정에서 민감정보 추론 방지 및 리스크 평가 기술 개발 (2026-2028, 참여연구자 (3책5공), Korea Internet \& Security Agency (한국인터넷진흥원))
    \item \small 토론과 소통으로 지식을 연결하는 AI 에이전트 기술 개발 (2026-2028, 참여연구자 (3책5공), Korea Creative Content Agency(한국콘텐츠진흥원))
    \item \small AI 보안 특화 'AI 에이전트' 개발 선행 연구 (2026,  참여연구자 (용역과제), 한국정보보호학회 산하 AI보안연구회)
    \item \small 국가·공공기관 가이드라인 준수를 위한 AI 시스템 보안 위협 자동화 기술점검 도구 개발 (2026, 참여연구자, KIST)
    \item \small AI 연구용 컴퓨팅 지원 프로젝트 (파라미터 수준의 정보 제어 원천기술 연구 및 파라미터 스튜디오 플랫폼 개발) (2026,  참여연구자, KETI)
    \item \small 지역지능화혁신인재양성 (ITRC) (2026-2029,  참여연구자, IITP)
\end{itemize}

\seclabel{Patents}
\textbf{Domestic Patents Applications}
\begin{itemize}
    \item \small{\textbf{인공지능 모델의 성능을 평가하는 벤치마크를 검증하는 전자 장치 및 이의 동작 방법} \\
    박찬준, \textbf{장지윤} \\
    \textit{Korean Patent Application, filed on August 7, 2026 (Application No. 10-2026-0147704).}}
\end{itemize}

\seclabel{Awards/Fellowships}
% No awards yet. Uncomment the block below to add them.
% \textbf{Awards}
% \begin{itemize}
%     \item \small{\placeholder{AWARD NAME}, \placeholder{AWARDING ORGANIZATION} (\placeholder{YEAR})}
% \end{itemize}
% \vspace{0.8em}
\textbf{Fellowships}
\begin{itemize}
    \item \small{Research Assistantship (RA) Scholarship, Soongsil University (2026)}
    \item \small{Merit-Based Scholarship (Departmental Recommendation), Soongsil University (2026)}
\end{itemize}

\seclabel{Teaching}
\textbf{Soongsil University}
\begin{itemize}
    \item \small{\textbf{Database}, Teaching Assistant (Spring 2026)}
\end{itemize}

\seclabel{Talks}
\textbf{Domestic Talks}
\begin{itemize}
    \item \small{\textbf{미소테크 (Misotech)}, API 및 Agent 실습 강연 (Hands-on Lecture on APIs and Agents), Jul 2026}
\end{itemize}

\end{document}
